\scanpage{56}
\margin{$\leftarrow$}Drapeaux d'un polyèdre convexe -- drapeaux d'un ens ordonné. Repères d'un polyèdre convexe

\uline{Prop.} Soit $C$ un polyèdre convexe, $\operatorname{Fac}(C)$ le polyèdre combinatoire, associé, $\operatorname{Rep}(C)=\operatorname{Rep}(\operatorname{Fac}(C))$ l'ens des repères. Alors les \deleted{applications} hom de gr
\[
\operatorname{Aut}_{\mathrm{aff}}(C)\xrightarrow{\alpha}\operatorname{Aut}_{\mathrm{ord}}(\operatorname{Fac}(C))\xrightarrow{\beta}\deleted{\unsure{S}}\operatorname{Aut}_{\mathrm{ens}}\operatorname{Rep}(C)
\]
\deleted{sont} injectives. Plus précisément pour $\beta$, $\operatorname{Aut}_{\mathrm{ord}}(\operatorname{Fac}(C))$ opère \uline{librement} sur $\operatorname{Rep}(C)$ (a fortiori, $\alpha$ l'étant aussi, $\operatorname{Aut}(C)$ opère librement)

En effet, $\alpha$ est inj. car un autom $u$ de $C$ est connu quand on connaît sa restriction aux pts extrémaux.

Soit $u$ un autom de $\Phi=\operatorname{Fac}(C)$, fixant un drapeau $(\unsure{x_0|F_0}\ldots x_{n-1},x_n)$. Par hyp de récurrence il \deleted{fixe $P_{x_{n-1}}$} est l'identité sur $\Phi_{x_{n-1}}$. D'autre part, s'il est l'identité sur un \deleted{facette} $P_{\deleted{\unsure{x_{n-1}}}x}$, $x$ une facette de dim $n-1$, et si $y$ est une facette telle que $x\wedge y$ est de dim $n-2$, alors il est l'identité sur $\Phi_y$ -- car il fixe un drapeau de $P_y$. Il faut voir que la rel. d'équiv. \inserted{sur $\unsure{\Phi_{n-1}}$} engendrée par ``$x$ et $y$ ont une facette de dim $n-2$ en commun'' est la relation grossière. \deleted{\unsure{Soit} en effet}

Argument topologique de connexité\ldots

\uline{Déf} $C$ \deleted{\unsure{est}} est dit polyèdre \uline{régulier} si $\operatorname{Aut}_{\mathrm{aff}}(C)$ est transitif (donc simplement transitif) sur $\operatorname{Rep}(C)$, homogène si pour tt $0\unsure{<|\leq}i\leq n\deleted{\unsure{+1|-1}}$, $\exists\,d_i\in\mathbb N$ tel que $\forall F\in\operatorname{Fac}_{i-1}(C)$, le cardinal de
% archive sheet number: 55
